% ============================================================
% Empirical Strategy chapter — Contract-Year Effect in Professional Football
% Selection-first framing; pairs with results_revised.tex.
% Self-contained: citations are written author-year in the prose,
% so the file compiles without a .bib. Convert to \citep/\citet later.
% ============================================================

\section{Empirical strategy}
\label{sec:strategy}

This chapter sets out how the data are used to learn about the contract-year
question, and why the design is built the way it is. Its organising decision is to
treat contract proximity not as a treatment that raises or lowers a player's
effort, but as a force that governs \emph{who is observed performing}. That choice
is not a stylistic preference; it is forced by the structure of the setting and of
the data. A performance rating exists only for a player his club selects to field,
and, as the analysis will show, that selection is exactly what contract proximity
moves. The strategy therefore separates two objects throughout: the fully observed
\emph{selection margin}---whether and how much a player is played---and the
selectively observed \emph{performance margin}---the per-minute rating, conditional
on appearing. The first is identified directly; the second, it will be argued, can
at best be bounded. The chapter proceeds from the identification philosophy and the
peculiar nature of the ``treatment'' in this setting, through the two-margin split
and the estimating equations, to the event study that is the study's main exhibit,
the partial-identification treatment of the performance margin, the regression
discontinuity that is tested and rejected, and finally inference, the robustness
protocol, and the principal threats to identification.

\subsection{Identification philosophy}
\label{sec:strategy-philosophy}

Three commitments discipline everything that follows. The first is a preference for
fully observed outcomes over selected, formula-mediated ones. Playing time is
recorded for every player-month in the panel; a rating is recorded only when the
player appears, and is then a composite of scoreboard events, minutes and team
performance rather than a direct reading of individual effort. Where the two
margins disagree, the design leans on the one that does not require a selection
correction to interpret. The second commitment is to identify from within-contract
time variation rather than from comparisons between players. Players who are deep
into a contract differ from players who are not in ways that are largely
unobservable, so the credible variation is the movement of a single player toward
his own contract's expiry, holding his identity fixed. The third commitment is to
audit every candidate effect to destruction---through season-by-season
re-estimation, a full specification grid, placebo cutoffs, event-time inspection,
and selection bounds---and to report a well-audited null in preference to a fragile
positive. These commitments explain why the chapter accumulates identification
strength in stages and why it treats the failure of a design as a result in its own
right.

\subsection{Contract expiry as an anticipated within-spell state}
\label{sec:strategy-treatment}

It is tempting to treat entry into the six-month Bosman window as a treatment and
to ask for its effect. The setting does not support that reading. Crossing into the
window is not an action any agent takes; it is a state reached mechanically as
calendar time advances within a fixed-term contract, and every player who remains
under contract reaches it. There is no take-up decision, no compliance margin, and
therefore none of the apparatus---instruments, local average treatment
effects---that a genuine take-up problem would call for. What varies is not whether
a player is treated but \emph{when}, along a spell whose clock runs
deterministically.

This places the setting in the literature on dynamically assigned, anticipated
contractual transitions rather than in the programme-evaluation literature on
treatment take-up. Abbring and van den Berg (2003) formalise the object closest to
ours: a treatment defined by the moment it is reached within a duration, identified
from the timing of events under a no-anticipation restriction rather than from a
binary treated-versus-control contrast. Two consequences follow for the design.
First, the legitimate comparison group is not players who never reach expiry but
the same player \emph{before} he reaches it. Defining ``controls'' as players who
never enter the window would condition on survival in the contract, because
whether and when a player reaches expiry is itself an outcome of the club's and the
player's earlier decisions; Fredriksson and Johansson (2008) show precisely that
such a static control definition, built on the future path of treatment, biases the
estimated effect, and that the correct comparison is the not-yet-treated (see also
Sianesi 2004). This is the formal justification for the within-spell design adopted
below. Second, because the Bosman date is perfectly anticipated and cannot be
manipulated, behaviour and allocation adjust smoothly as it approaches rather than
jumping when it is crossed. That is an anticipation problem, and it is distinct
from the manipulation problem that motivates the usual density diagnostics
(McCrary 2008); it is also the reason, developed in
Section~\ref{sec:strategy-rdd}, that a discontinuity design at the threshold is not
merely underpowered here but mis-specified.

\subsection{Two margins, two identification problems}
\label{sec:strategy-margins}

The outcome structure splits the analysis into two problems that must be kept
separate. The extensive, or selection, margin---whether a player plays at all in a
month, and the minutes and matches he accumulates---is observed for all 152{,}817
player-months from the Transfermarkt playing records, independently of whether he
appears. It is therefore identified directly: no correction is needed to interpret
a within-player, within-spell change in playing probability as the club's
allocation response. The intensive, or performance, margin---the monthly match
rating---exists only when the player is selected to appear, and is missing exactly
for the player-months in which he is not. This is a selective-observation problem
of the classic kind. It is the same structure that arises when wages are observed
only for those who work, so that a regression on the observed sample confounds the
outcome equation with the equation governing who is observed (Heckman 1979; Gronau
1974); ratings are to latent performance what observed wages are to the wage-offer
function, and playing is the analogue of employment.

What makes the performance margin especially hard here is that the selection runs
on the outcome itself. Clubs field the players they currently judge most useful, so
the set of appearances is selected on close correlates of the very quantity the
rating is meant to measure---a sorting-on-outcomes structure of the Roy (1951)
type, which is the least favourable case for recovering the latent distribution
from the observed one. The direction of the resulting bias is signable even when
its magnitude is not. Because expiring players who are withdrawn are the ones the
club rates least, the players who continue to appear near expiry are positively
selected, so any conditional rating comparison is biased \emph{upward} relative to
a true negative effect. A conditional null on the rating is thus consistent with a
negative, a zero, or a positive population effect; conditioning on selection cannot
tell them apart. This is why the strategy does not attempt to point-identify the
performance margin and instead bounds it (Section~\ref{sec:strategy-bounds}), and
why the weight of the analysis rests on the selection margin, which it can identify.

\subsection{Estimating equations and the specification ladder}
\label{sec:strategy-ladder}

The analysis climbs a ladder of specifications of increasing identification
strength, reporting the whole ladder rather than a single preferred cell so that
the reader can see where each candidate effect survives or dies. The base outcome
model is a linear two-way fixed-effects regression of the form
\[
  y_{it} \;=\; \beta\,\mathrm{Bosman}_{it} \;+\; \alpha_i \;+\; \gamma_t \;+\; \varepsilon_{it},
\]
where $y_{it}$ is the outcome for player $i$ in month $t$, $\mathrm{Bosman}_{it}$
indicates that the player's contract has six months or fewer remaining, $\alpha_i$
is a player fixed effect, $\gamma_t$ a time fixed effect, and standard errors are
clustered by player. The Bosman indicator is measured contemporaneously from
contract status, using no information about the player's future path, so the
regressor cannot leak future outcomes into the present.

The ladder tightens this in four steps. The first replaces the calendar month
fixed effect with a league-by-month fixed effect, so that a player is compared only
to others in the same league in the same month; this absorbs the shifting league
composition of the panel, which is the design's largest single threat. The second
replaces the player fixed effect with a player-by-contract-spell fixed effect. This
is the decisive step: within a single spell, time-to-expiry advances mechanically
and cannot be a state the player has sorted into, so the comparison is of one
player with himself along one contract. Identification here is real rather than
nominal, because a majority of players are observed across more than one spell
(2{,}082 of 5{,}412 players, spanning 7{,}901 spells). The third step combines
spell fixed effects with league-by-month fixed effects, the strictest general
specification. The fourth uses club-by-month fixed effects, comparing an expiring
player directly with his own teammates in the same month and thereby differencing
out anything common to the club's situation---form, injuries, managerial change.
Alongside the single Bosman indicator, the same models are estimated with six-month
distance-to-expiry bins, which trace the profile across the contract without
imposing that the entire response occurs at the legal threshold.

\subsection{The contract-cycle event study}
\label{sec:strategy-eventstudy}

The single Bosman indicator compresses the contract into two states; the study's
main exhibit instead estimates the whole trajectory. Each outcome is regressed on a
full set of months-to-expiry indicators, from twenty-four-or-more down to zero with
eighteen months as the omitted reference, within player-by-spell and
league-by-month fixed effects. This traces the selection rule and the conditional
rating continuously across the contract and lets the data, rather than the
six-month convention, say where any response occurs. It is also the natural design
given the anticipation structure of Section~\ref{sec:strategy-treatment}: an
anticipated, non-manipulable deadline predicts a smooth, forward-looking adjustment,
which an event study can display and a two-state indicator cannot.

Two methodological cautions attach to this design and are stated openly. First,
every player in the panel is eventually treated in the sense that every contract
runs toward expiry, so there is no never-treated comparison group. In staggered
settings of this kind the conventional two-way fixed-effects estimator can place
negative weights on some comparisons, effectively using already-treated units as
controls when treatment effects are heterogeneous over the cycle (Goodman-Bacon
2021; de Chaisemartin and D'Haultf\oe uille 2020; Sun and Abraham 2021). The
within-spell event study mitigates this by reading effects off the event-time path
relative to a fixed within-spell reference rather than off a pooled dummy, and the
heterogeneity-robust estimators developed for anticipated, all-eventually-treated
designs---in particular the group-time approach of Callaway and Sant'Anna (2021),
which admits an explicit limited-anticipation window---are the formal home for this
exhibit and the natural route for a fuller treatment of the renewal margin. Second,
because renewal is itself selected, the event study around the first observed
contract extension is interpreted with the same caution as the rating regressions:
the interesting quantity is the within-player path around signing, not the level
difference between renewed and non-renewed players.

\subsection{Partial identification of the performance margin}
\label{sec:strategy-bounds}

The selective-observation problem of Section~\ref{sec:strategy-margins} means the
conditional rating effect is not point-identified, and the strategy treats this as
a feature to be characterised rather than a defect to be assumed away. The
selection margin is estimated directly on the full panel, where playing is always
observed, and delivers the object the design can identify without correction. For
the performance margin the analysis reports bounds. Following Lee (2009), the
group that is over-observed near expiry---the players who continue to appear---is
trimmed by its differential observation share within league-month cells, under the
monotonicity assumption that contract proximity moves selection in one direction,
yielding sharp bounds on the Bosman rating gap. This is the appropriate posture for
a Roy-type selection on the outcome, where worst-case reasoning rather than a
point estimate is what the data support (Manski's partial-identification
programme; Mourifi\'e, Henry and M\'eango 2018 give sharp bounds for exactly a Roy
model with selection on potential outcomes). The result to be reported in the next
chapter is that the identified set spans zero: given the selection into playing, the
conditional rating cannot even sign a performance effect. Separately, the
possibility that reaching expiry is itself a selected state---because clubs may
already have decided not to renew, or because form or fitness has already
declined---is the dynamic-selection analogue of the Ashenfelter dip (Heckman and
Smith 1999); the within-spell design addresses selection \emph{into the
trajectory}, while the bounds address selection \emph{into appearing}, and both are
required because they are different problems.

\subsection{The regression discontinuity, and why it is rejected}
\label{sec:strategy-rdd}

A regression discontinuity at 180 days to expiry was implemented as the design that
would, in principle, offer the cleanest local comparison, estimating a local linear
model in days-to-expiry around the Bosman cutoff with month fixed effects and
clustering by player. It is reported, and then rejected, as a negative exhibit. The
running variable is measured in monthly steps, so support within a narrow bandwidth
is almost absent---of order a hundred observations against nine hundred at a
thirty-day bandwidth---and the local-linear estimate is implausibly large and of
the wrong sign relative to the raw gap. Placebo cutoffs at 150 and 210 days fire at
conventional significance, confirming that the design is detecting local slope and
composition rather than anything at the legal threshold. The correct reading of this
failure is the one anticipated in Section~\ref{sec:strategy-treatment}: because free
agency at six months is perfectly foreseen and cannot be manipulated, there is no
discontinuous change in the club's selection problem at 180 days for a discontinuity
design to exploit. The failure is thus a prediction of the identification analysis,
not an accident of the sample, and it is precisely what motivates preferring the
smooth event study to the discontinuity.

\subsection{Inference}
\label{sec:strategy-inference}

Inference recognises that observations are clustered within players and within
time. The baseline clusters standard errors by player; robustness specifications
cluster two-way by player and month, and, in the teammate comparisons, by club.
Because the heterogeneity analysis inspects roughly a dozen subgroups per
hypothesis, isolated significant cells are expected under the null, and the analysis
reports Benjamini-Hochberg adjusted $q$-values across subgroups within each
model-outcome-term family; it also notes explicitly that this adjustment does not
repair a season-stability failure, which is a different problem from multiplicity.
Finally, every null is accompanied by a minimum detectable effect, so that an
absence of significance is shown to be informative rather than merely uncertain:
for the Bosman rating coefficient the minimum detectable effect at eighty per cent
power is about 0.06--0.07 standard deviations, so effects larger than this are ruled
out with good power.

\subsection{Robustness protocol}
\label{sec:strategy-robustness}

Candidate effects are subjected to a fixed protocol designed to break them. Every
pooled effect is re-estimated season by season, and any effect that is significant
in the pooled data but vanishes within each season is treated as a composition
artefact of the panel's changing league mix rather than as a behavioural response.
The full specification grid is reported rather than a single favoured cell, so that
sign instability across fixed-effects structures is visible. Threshold claims are
checked against placebo cutoffs. Any dynamic claim is checked against its event-time
path, so that a pooled coefficient must correspond to a coherent trajectory to be
believed. And all time-varying regressors are assigned without look-ahead: contract
status is contemporaneous, and club financial variables are lagged to the financial
year ending before each season, strictly backward-looking; an earlier result that
depended on improperly time-averaged financial assignment did not survive correct
lagging and is retained only as a cautionary note.

\subsection{Threats to identification}
\label{sec:strategy-threats}

Three threats deserve to be named at the outset. The first and largest is that the
panel's league composition shifts substantially across seasons, from the top five
leagues only in the earliest seasons to some forty leagues later; this is the
motivation for league-by-month fixed effects and for the season-stability audit, and
it remains the single largest caveat on any pooled estimate. The second is that part
of the selection-margin decline may be mechanical rather than discretionary---driven
by the timing of mid-season transfers, by injuries that run down contracts, or by
clubs protecting the sale value of a departing player---and separating these channels
would require transfer-event data that the panel does not contain; this is carried
forward to the discussion as a limitation on the interpretation, though not on the
existence, of the effect. The third is that the performance margin is only ever
observed for the selected, so that even a well-powered conditional null is a bounded
rather than a point statement, a limitation the strategy confronts directly rather
than assumes away. How football's anticipated, non-manipulable deadline maps to
other fixed-term labour markets---and to the closely analogous political
term-limits setting, where the same discipline-versus-selection problem has been
studied (Besley and Case 1995; List and Sturm 2006; Alt, Bueno de Mesquita and Rose
2011)---is taken up in the discussion.
